\chapter{Introduction}
\label{ch:intro}

\section{Motivation: the cost of deep learning in sensing systems}
\label{sec:intro:motivation}

Deep neural networks have become the default tool for perception tasks such as
object detection in video and land-cover classification from satellite imagery.
Their accuracy comes at a cost that is paid at every inference: the input must
be acquired, moved, stored and processed, and the network must be executed on
some hardware platform. In deployed sensing systems, a road-side traffic
camera, a drone, or a country-scale crop-mapping service, this cost is repeated
for every frame, every field and every season.

Most work on efficient deep learning targets the \emph{execution} of a given
network: faster software frameworks, compilers, reduced numerical precision,
pruning and specialised hardware. These techniques are essential, but they share
an implicit assumption: that the full input must be processed. This thesis
questions that assumption. In many sensing applications the input is highly
redundant with respect to the task, and the cheapest input is the one that is
never acquired.

\section{The big picture}
\label{sec:intro:bigpicture}

\paragraph{Where the cost goes.}
A deployed sensing system runs a pipeline (\cref{fig:intro:bigpicture}): a
sensor or an archive delivers the input, the input is acquired and transferred,
stored and pre-processed, and finally a deep neural network turns it into a
prediction. Only the last stage is what is usually called inference, but every
stage costs time, energy, bandwidth or money, and every stage is paid again for
each frame, field and season. The balance between the stages depends on scale.
For a single camera attached to an embedded GPU, the network execution
dominates. For a crop map of a whole country, the network is cheap compared with
retrieving, storing and pre-processing tens of multispectral acquisitions per
field. As sensing systems scale, the dominant cost moves from the processor to
the data.

\paragraph{Thesis statement.}
This thesis argues that the most effective way to make deep learning inference
in sensing systems efficient is to process less input, not only to process it
faster: the part of the input a model actually needs can be identified
cheaply, before or during inference, and the rest need not be acquired, stored
or processed. The question therefore shifts from \emph{how} to execute a
network efficiently to \emph{what} part of the input needs to be processed at
all.

\begin{figure}[t]
  \centering
  \begin{tikzpicture}[
  font=\small,
  stage/.style={draw, rounded corners=2pt, align=center, text width=19mm,
                minimum height=11mm, fill=black!5},
  box/.style={draw, rounded corners=2pt, align=center, font=\footnotesize,
              inner sep=4pt},
  flow/.style={-{Stealth[length=2mm]}, thick},
  ctrl/.style={-{Stealth[length=1.8mm]}},
  fb/.style={-{Stealth[length=1.8mm]}, densely dashed},
  node distance=6mm and 5mm]

  \node[stage] (sel) {Input\\selection};
  \node[stage, right=of sel] (acq) {Acquire \&\\transfer};
  \node[stage, right=of acq] (sto) {Store \&\\pre-process};
  \node[stage, right=of sto] (inf) {DNN\\inference};
  \node[align=center, font=\footnotesize, left=4mm of sel] (src) {Sensor\\or\\archive};
  \draw[flow] (src) -- (sel);
  \draw[flow] (sel) -- (acq);
  \draw[flow] (acq) -- (sto);
  \draw[flow] (sto) -- (inf);

  \draw[fb] (inf.north) -- ++(0,7mm) -| (sel.north)
    node[pos=0.25, above, font=\footnotesize]
    {feedback at run time: prior detections (Ch.~4), confidence (Ch.~6)};

  \node[box, anchor=north, text width=63mm] (what)
    at ($(acq.south)+(0,-13mm)$)
    {\textbf{What} to process: reduces the data cost\\[2pt]
     \textbf{Where}: image regions (Ch.~4)\quad \textbf{Which}: bands (Ch.~5)\\
     \textbf{When}: acquisitions (Ch.~6)\quad \textbf{Which+When} (Ch.~7)};
  \node[box, anchor=north, text width=19mm] (how)
    at ($(inf.south)+(0,-13mm)$)
    {\textbf{How} to execute: reduces the compute cost\\(Ch.~3)};
  \draw[ctrl] (what.north -| sel.south) -- (sel.south);
  \draw[ctrl] (how.north) -- (inf.south);
\end{tikzpicture}

  \caption{The big picture of this thesis. A sensing pipeline pays a data cost
  (acquire, transfer, store, pre-process) and a compute cost (inference).
  \Cref{ch:how} reduces the compute cost by choosing how a fixed network is
  executed. \Cref{ch:where,ch:which,ch:when,ch:joint} add an input-selection
  stage that decides what enters the pipeline: which image regions, spectral
  bands and acquisitions. The selection is fixed at design time from training
  data or model attributions, or adapted at run time from the model's own
  output (dashed).}
  \label{fig:intro:bigpicture}
\end{figure}

\paragraph{Three axes of input redundancy.}
A sensing input has a spatial, a spectral and a temporal extent, and each is
redundant with respect to the task:
\begin{itemize}
  \item \textbf{Where} (space): most pixels of a traffic-camera frame contain no
  object, and objects move little between consecutive frames (\cref{ch:where}).
  \item \textbf{Which} (spectrum): neighbouring spectral bands of a
  multispectral image carry overlapping information, and a crop classifier needs
  only a few of them (\cref{ch:which}).
  \item \textbf{When} (time): after the discriminative part of the growing season
  has been observed, further acquisitions add little (\cref{ch:when}).
\end{itemize}
The axes are not independent: the bands a classifier needs may depend on the
dates it sees. \Cref{ch:joint} therefore selects bands and dates jointly, under
one data-volume budget.

\paragraph{One recipe, three axes.}
The chapters on input reduction follow the same recipe.
\begin{enumerate}
  \item \emph{Bound the redundancy.} Measure how much of the input can be removed
  at all, with a reference that is too expensive to deploy: the full-frame
  detector, an exhaustive search over band subsets, or a label-aware oracle over
  acquisitions.
  \item \emph{Find a cheap selector.} Replace the reference with a signal that is
  available at little cost: detections from previous frames, feature
  attributions of a trained model, statistics of the training set, or the
  classifier's own confidence. A selector is only useful if it costs much less
  than the input it saves.
  \item \emph{Test whether it transfers.} Check the selection on data it was not
  derived from: other datasets, regions, architectures and years.
\end{enumerate}
The selectors differ in when they decide. A \emph{static} selection is fixed at
design time and is the same for every input, such as a band set or an
acquisition calendar; it is simple to deploy and lets the data provider deliver
only what is needed. An \emph{adaptive} selection is made at run time for each
input, such as crop regions from prior detections or a stopping rule; it can
save more, but it must be computed online and its decisions depend on the
model being right. A recurring finding of this thesis is that good static
selections are hard to beat (\cref{ch:when}).

\paragraph{The thesis at a glance.}
\Cref{tab:intro:glance} summarises the chapters along these lines.

\begin{table}[t]
  \centering
  \footnotesize
  \caption{The thesis at a glance: what each chapter reduces, how the selection
  is made, and the main result.}
  \label{tab:intro:glance}
  \begin{tabularx}{\textwidth}{@{}l l X X X@{}}
    \toprule
    Ch. & Axis & Reduced & Selection signal (when) & Main result \\
    \midrule
    \ref{ch:how}   & How       & execution cost      & framework, platform, precision (design time)
                   & no single framework is best; latency and throughput optima differ \\
    \ref{ch:where} & Where     & image regions       & prior detections (run time, per frame)
                   & up to $1.7\times$ faster detection at maintained accuracy \\
    \ref{ch:which} & Which     & spectral bands      & model attributions (design time, per region)
                   & 3 bands retain 99.1\% of 10-band performance \\
    \ref{ch:when}  & When      & acquisitions        & searched calendar (design time); stopping head (run time)
                   & within 1 point of the full series with a fraction of the acquisitions \\
    \ref{ch:joint} & Which+When & band--date pairs   & joint search under a data-volume budget (design time)
                   & first results \\
    \bottomrule
  \end{tabularx}
\end{table}

\paragraph{Scope.}
The thesis spans two domains on purpose. Object detection in video on embedded
hardware (\cref{ch:how,ch:where}) is the setting where compute dominates;
crop classification from Sentinel-2 time series (\cref{ch:which,ch:when,ch:joint})
is the setting where data dominates. Seeing the same question in both shows that
input reduction is not a trick of one application. The thesis studies inference
only: models are trained once on full inputs, or on random subsets of them, and
the savings are measured at deployment.

\section{Research questions}
\label{sec:intro:rq}

\begin{description}
  \item[RQ1] How strongly do the choices of deep learning software framework and
  hardware platform affect inference latency and throughput, and is there a
  choice that is best across deployment scenarios? (\cref{ch:how})
  \item[RQ2] Can prior object locations in video be exploited to reduce the
  computation of a single-shot object detector while maintaining its accuracy?
  (\cref{ch:where})
  \item[RQ3] Can model feature attributions identify a minimal set of spectral
  bands for crop classification, close to the optimum found by exhaustive
  search, at a fraction of its cost? (\cref{ch:which})
  \item[RQ4] How few Sentinel-2 acquisitions does each field need for
  near-full-series accuracy, and can they be selected, and the retrieval
  stopped, without labels at inference time? (\cref{ch:when})
  \item[RQ5] Do spectral and temporal input reductions compose, and how should
  they be selected jointly under a common data-volume budget? (\cref{ch:joint})
\end{description}

\section{Contributions}
\label{sec:intro:contrib}

\begin{enumerate}
  \item A systematic evaluation of inference performance across deep learning
  frameworks, hardware platforms, batch sizes and precisions, showing that no
  single framework is optimal and that the optimum differs for latency and
  throughput (\cref{ch:how}).
  \item A crop-based detection method that turns prior object locations into
  low-resolution crop regions, with a crop-size limit as an explicit
  speed/accuracy knob; up to $1.7\times$ (UA-DETRAC) and $1.6\times$ (UAVDT)
  speed-up over SSD-ML (\cref{ch:where}).
  \item An attribution-ensemble method for automatic spectral band selection,
  validated against 3{,}675 exhaustive trainings across six Turkish regions
  and three architectures; three bands retain 99.1\% of 10-band performance
  while reducing storage by 70--90\% (\cref{ch:which}).
  \item A retrieval-efficient acquisition framework, consisting of an any-subset
  classifier, a per-region searched acquisition calendar, and a label-free
  learned stopping head, that reaches near-full-series accuracy with a small
  fraction of the acquisitions, and composes with early-classification methods
  (\cref{ch:when}).
  \item Evidence that label-aware oracles, commonly used as teachers for
  acquisition policies, mostly reflect selection bias, which explains why
  learned per-field policies fail to beat good static schedules
  (\cref{ch:when}).
  \item A joint spectral--temporal selection under a common data-volume cost,
  showing how the two reductions compose and where they fail to transfer
  (\cref{ch:joint}).
\end{enumerate}

\section{Thesis outline}
\label{sec:intro:outline}

\Cref{ch:background} introduces the background and a common cost--accuracy
framework that is used throughout the thesis. \Cref{ch:how} studies how the
execution of a fixed network depends on the software and hardware, and argues
that larger gains require reducing the input. \Cref{ch:where,ch:which,ch:when}
reduce the input along space, spectrum and time, respectively.
\Cref{ch:joint} combines the spectral and temporal reductions. \Cref{ch:conclusion}
answers the research questions and discusses future directions, including
on-board processing of satellite data, which brings the embedded perspective of
the first chapters back to the remote-sensing setting of the last.

\section{Publications}
\label{sec:intro:pubs}
The chapters of this thesis are based on the following publications:
\cref{ch:how}~\cite{ulker2020reviewing}, \cref{ch:where}~\cite{ulker2022accelerating},
\cref{ch:which}~\cite{ulker2026bandselection}, and
\cref{ch:when,ch:joint}~\cite{ulker2027acquisitions}.
